\documentclass[a4paper,11pt]{article}
\usepackage[T1]{fontenc}
\usepackage[utf8]{inputenc}
\usepackage[french]{babel}
\usepackage{lmodern}
\usepackage{xcolor}
\usepackage[margin=1.4cm]{geometry}
\usepackage{tikz}
\pagestyle{empty}
\usetikzlibrary{arrows.meta,positioning}

\definecolor{astblue}{HTML}{DCEEFF}
\definecolor{astgreen}{HTML}{DFF5E1}
\definecolor{astyellow}{HTML}{FFF3CC}
\definecolor{astpink}{HTML}{FBE1EA}
\definecolor{astgray}{HTML}{F4F4F4}
\definecolor{astline}{HTML}{5A6B7A}

\tikzset{
  box/.style={
    draw=astline,
    rounded corners=3pt,
    thick,
    align=left,
    font=\sffamily\footnotesize,
    inner sep=6pt
  },
  title/.style={
    font=\sffamily\bfseries\Large,
    text=black
  },
  subtitle/.style={
    font=\sffamily\bfseries\large,
    text=black
  },
  arrow/.style={
    -{Latex[length=3mm,width=2mm]},
    thick,
    draw=astline
  },
  note/.style={
    font=\sffamily\footnotesize,
    text=black!75,
    align=left
  }
}

\begin{document}

\begin{center}
{\sffamily\bfseries\LARGE Kairos IR -- Vue d'ensemble}

\vspace{2mm}
{\sffamily\small Les IR commencent après le \texttt{Normalized\_program}. On distingue l'IR des obligations de preuve et le kernel exporté.}
\end{center}

\vspace{4mm}

\begin{center}
\begin{tikzpicture}[x=0.84cm,y=0.92cm]

\node[subtitle] at (8.2, 11.2) {Page 1 -- Proof\_obligation\_ir};

\node[box, fill=astgray, text width=4.0cm, anchor=north west] (input) at (0.0, 10.2) {
  \textbf{Entrée}\\
  \rule{\linewidth}{0.4pt}\\
  \texttt{Normalized\_program.node}\\
  programme normalisé\\
  + contrats de transition générés\\
  (\texttt{requires}/\texttt{ensures})
};

\node[box, fill=astblue, text width=4.2cm, anchor=north west] (raw) at (4.7, 10.2) {
  \textbf{raw\_node}\\
  \rule{\linewidth}{0.4pt}\\
  \texttt{node\_name, inputs, outputs, locals}\\
  \texttt{control\_states, init\_state}\\
  \texttt{instances, pre\_k\_map}\\
  \texttt{transitions : raw\_transition list}\\
  \texttt{assumes, guarantees}\\[1mm]
  \textit{raw\_transition}\\
  \texttt{src\_state, dst\_state}\\
  \texttt{guard, guard\_iexpr}\\
  \texttt{body\_stmts}
};

\node[box, fill=astgreen, text width=4.4cm, anchor=north west] (annotated) at (9.7, 10.2) {
  \textbf{annotated\_node}\\
  \rule{\linewidth}{0.4pt}\\
  \texttt{raw : raw\_node}\\
  \texttt{transitions : annotated\_transition list}\\
  \texttt{coherency\_goals}\\
  \texttt{user\_invariants}\\
  \texttt{state\_invariants}\\[1mm]
  \textit{annotated\_transition}\\
  \texttt{raw : raw\_transition}\\
  \texttt{requires : contract\_formula list}\\
  \texttt{ensures : contract\_formula list}
};

\node[box, fill=astyellow, text width=4.4cm, anchor=north west] (verified) at (14.8, 10.2) {
  \textbf{verified\_node}\\
  \rule{\linewidth}{0.4pt}\\
  \texttt{node\_name, inputs, outputs, locals}\\
  \texttt{control\_states, init\_state}\\
  \texttt{instances}\\
  \texttt{transitions : verified\_transition list}\\
  \texttt{assumes, guarantees}\\
  \texttt{coherency\_goals}\\
  \texttt{user\_invariants, state\_invariants}\\[1mm]
  \textit{verified\_transition}\\
  \texttt{src\_state, dst\_state}\\
  \texttt{guard, guard\_iexpr}\\
  \texttt{body\_stmts}\\
  \texttt{pre\_k\_updates}\\
  \texttt{requires / ensures}
};

\node[box, fill=astpink, text width=5.3cm, anchor=north west] (passes) at (5.0, 4.1) {
  \textbf{Passes}\\
  \rule{\linewidth}{0.4pt}\\
  \texttt{raw\_obligation\_generation}\\
  \texttt{triple\_annotation}\\
  \texttt{history\_lowering}
};

\node[box, fill=astgray, text width=7.1cm, anchor=north west] (meaning) at (11.4, 4.1) {
  \textbf{Sens de cette phase}\\
  \rule{\linewidth}{0.4pt}\\
  1. matérialiser la structure exécutable du programme\\
  2. attacher les contrats \texttt{requires}/\texttt{ensures}\\
  3. expliciter l'historique \texttt{pre\_k} dans une forme stabilisée\\[1mm]
  \texttt{Proof\_obligation\_ir} reste centré sur les\\
  transitions du programme et leurs obligations.
};

\node[note, text width=17.6cm, anchor=north west] (legend1) at (0.0, 0.7) {
  \textbf{Lecture rapide}\\
  \texttt{raw\_node} = structure exécutable sans triples.\\
  \texttt{annotated\_node} = mêmes transitions, enrichies par les contrats calculés sur le programme normalisé.\\
  \texttt{verified\_node} = version où l'historique a été abaissé et où les obligations sont prêtes à être consommées par le kernel.
};

\draw[arrow] (input.east) -- (raw.west);
\draw[arrow] (raw.east) -- (annotated.west);
\draw[arrow] (annotated.east) -- (verified.west);
\draw[arrow] (annotated.south) -- (passes.north);
\draw[arrow] (verified.south) -- (meaning.north);
\draw[arrow] (passes.east) -- (meaning.west);

\end{tikzpicture}
\end{center}

\newpage

\begin{center}
\begin{tikzpicture}[x=0.84cm,y=0.92cm]

\node[subtitle] at (8.2, 11.2) {Page 2 -- Proof\_kernel\_ir};

\node[box, fill=astgray, text width=4.5cm, anchor=north west] (input2) at (0.0, 10.1) {
  \textbf{Entrée}\\
  \rule{\linewidth}{0.4pt}\\
  \texttt{verified\_node}\\
  + automates \texttt{assume} / \texttt{guarantee}\\
  + analyse produit / appels\\
  selon la chaîne kernel
};

\node[box, fill=astgreen, text width=4.8cm, anchor=north west] (kp) at (5.1, 10.1) {
  \textbf{Passes proof\_kernel}\\
  \rule{\linewidth}{0.4pt}\\
  \texttt{proof\_kernel\_program}\\
  \texttt{proof\_kernel\_clauses}\\
  \texttt{proof\_kernel\_step\_contracts}\\
  \texttt{proof\_kernel\_calls}\\
  \texttt{proof\_kernel\_ir}
};

\node[box, fill=astblue, text width=7.3cm, anchor=north west] (kernel) at (10.8, 10.1) {
  \textbf{Proof\_kernel\_ir.node\_ir}\\
  \rule{\linewidth}{0.4pt}\\
  \texttt{reactive\_program}\\
  \texttt{assume\_automaton}\\
  \texttt{guarantee\_automaton}\\
  \texttt{initial\_product\_state}\\
  \texttt{product\_states}\\
  \texttt{product\_steps}\\
  \texttt{generated\_clauses}\\
  \texttt{symbolic\_generated\_clauses}\\
  \texttt{proof\_step\_contracts}\\
  \texttt{instance\_relations}\\
  \texttt{call\_site\_instantiations}\\
  \texttt{signature / exported summaries}
};

\node[box, fill=astyellow, text width=6.3cm, anchor=north west] (meaning2) at (2.2, 4.3) {
  \textbf{Sens de cette phase}\\
  \rule{\linewidth}{0.4pt}\\
  Le kernel n'est plus seulement l'IR des transitions.\\
  Il assemble :\\
  - programme réactif\\
  - automates de spécification\\
  - produit explicite\\
  - clauses générées\\
  - contrats par transition du produit\\
  - artefacts exportables
};

\node[box, fill=astgray, text width=5.7cm, anchor=north west] (outputs) at (10.8, 4.3) {
  \textbf{Sorties}\\
  \rule{\linewidth}{0.4pt}\\
  \texttt{.kobj}\\
  rendus / debug IR\\
  traduction Why3\\[1mm]
  Le backend Why ne doit faire\\
  qu'une traduction structurelle.
};

\node[note, text width=17.6cm, anchor=north west] (legend2) at (0.0, 0.7) {
  \textbf{Lecture rapide}\\
  \texttt{Proof\_kernel\_ir} est le vrai noyau exporté de la chaîne.\\
  On y trouve le produit programme $\times$ automates, les clauses qui en découlent, puis les contrats attachés aux transitions du produit.\\
  C'est cette structure qui sert de source de vérité pour le \texttt{.kobj} et pour Why3.
};

\draw[arrow] (input2.east) -- (kp.west);
\draw[arrow] (kp.east) -- (kernel.west);
\draw[arrow] (kp.south) -- (meaning2.north);
\draw[arrow] (kernel.south) -- (outputs.north);
\draw[arrow] (meaning2.east) -- (outputs.west);

\end{tikzpicture}
\end{center}

\end{document}
